\section{Gurobi Fixed-Interval Implementation}
The stable inner backend is a sparse tailored fixed-interval MILP solved by
Gurobi's native branch-and-cut. Gurobi manages presolve, LP relaxations,
branching, node selection, incumbent processing, best bounds, and termination.
Official deterministic settings are \texttt{Threads=1}, \texttt{Seed=0},
\texttt{Presolve=Auto}, \texttt{MIPGap=0}, and \texttt{MIPGapAbs=0}.
Branching priorities are absent, PreCrush remains at its default, and no
MIPNODE user-cut callback is registered.

The fixed-interval model is freshly constructed for each ledger leaf. Its
canonical fingerprint binds variables, rows, domains, objective, input, solver
settings, and executable identity. Bounds are accepted only when model scope,
lifecycle, and fingerprint checks pass. Integer candidates are reconstructed
and checked against original route, load, inventory, Gini, penalty, and
objective definitions before they may update the incumbent.

Round 54's inventory--route experiment was intentionally outside this stable
path. It solved a fresh root LP, separated the most violated directed cut by an
exact deterministic maximum-closure/min-cut construction with the depot fixed
outside the station subset, added canonical nonduplicate rows, and rebuilt a
fresh terminal MIP with integer types restored. Root closure had no hidden
node, size, instance, or time dispatch. Invalid closure triggered F0 fallback.
IR1 used the two mixed inequalities to full closure; IR2 added projected
companions; IR3 was a predeclared one-pass option. Only IR1 entered the live
comparison, and its negative result prevented integration into the stable
preset.

Round 55 implements three additional default-off fixed-interval policies.
MC4 adds the four aggregate McCormick inequalities for $Z_i=G Y_i$ on each
station domain. VD--P replaces the product-only bit block by exact
station-state selectors and perspective variables. VD--J jointly represents
the product, ratio, and penalty on the same selectors while retaining the
original absolute-deviation inequalities. These policies use the same native
branch-and-cut settings and register no callback. Mathematical exactness and
root strength are audited separately from total MIP performance.

Plain Gurobi is retained as a contextual benchmark. Its bound or certificate
never enters the tailored interval ledger. Historical CPLEX-managed callback
experiments are likewise contextual and are not the current implementation.
